\section{Заключение}
\label{sec:conclusion}
\porozhdeno{\Claim{C01}{При фиксированных ASV и алгоритме выбор готовой модели меняет покрытие назначениями, таксономические линии и таксономические профили исследованных почв. В дополнительных профилях на уровне семейств изменяются и множество учитываемых ридов, и трактовка терминальной метки uncultured. Поэтому сопоставление требует одновременно указывать покрытие, знаменатель и правило для неназначенной массы, отличая пустой отбор от неопределённого семейства. Полученные различия характеризуют чувствительность сохранённых назначений; они не определяют наиболее точную модель и не устанавливают экологические причины состава почв.}}
\napisano{\Claim{C02}{Для воспроизводимости сравнения необходимо документировать версии референсных баз и готовых моделей, параметры классификации и правила формирования анализируемых наборов. Более строгое разграничение влияния референсной базы и условий подготовки классификатора требует сравнения моделей, обученных по единому протоколу. Следующим этапом может стать проверка на образцах с известным таксономическим составом, позволяющая оценить не только полноту, но и корректность назначений.}}
